\subsection{So sánh các chiến lược bổ sung dữ liệu thiếu (Imputation)}

\subsubsection{Hồi quy}

\tablename~\ref{tab:regression-sens-imputer} trình bày độ nhạy của các mô hình
hồi quy trước các tỉ lệ dữ liệu thiếu khác nhau khi áp dụng các chiến lược bổ
sung dữ liệu (mean, median và kNN).

\import{\tablespath}{regression-sens-imputer}

Trước hết, có thể nhận thấy một xu hướng nhất quán rằng khi tỉ lệ dữ liệu thiếu
tăng từ $0.1$ lên $0.3$, độ nhạy của tất cả các mô hình đều tăng đáng kể trên
mọi chỉ số (MSE, RMSE, MAE và $R^2$). Điều này cho thấy việc bổ sung dữ liệu
không thể hoàn toàn bù đắp thông tin bị mất, đặc biệt trong trường hợp dữ liệu
thiếu ở mức cao.

Xét theo từng phương pháp, kNN imputation cho kết quả tốt nhất một cách nhất
quán trên hầu hết các mô hình và mức thiếu dữ liệu. Cụ thể, các giá trị độ nhạy
của kNN luôn thấp hơn đáng kể so với mean và median, ví dụ với OLS tại mức
thiếu $0.1$, MSE sensitivity giảm từ $0.2733$ (mean) và $0.2967$ (median) xuống
còn $0.1796$. Xu hướng này tiếp tục được duy trì ở các mức thiếu cao hơn, cho
thấy kNN có khả năng bảo toàn cấu trúc cục bộ của dữ liệu tốt hơn so với các
phương pháp dựa trên thống kê toàn cục.

Mean imputation đóng vai trò như một phương pháp trung gian, thường cho kết quả
tốt hơn median nhưng vẫn kém hơn kNN. Phương pháp này có ưu điểm về tính đơn
giản và chi phí tính toán thấp, tuy nhiên việc thay thế bằng giá trị trung bình
làm giảm độ biến thiên của dữ liệu, từ đó ảnh hưởng đến khả năng học của mô
hình.

Ngược lại, median imputation thường cho kết quả kém nhất trong các trường hợp
thử nghiệm. Độ nhạy cao hơn cho thấy phương pháp này không phù hợp khi dữ liệu
không bị lệch hoặc khi mối quan hệ giữa các biến mang tính phi tuyến.

Xét theo từng nhóm mô hình, các mô hình tuyến tính (OLS, MBGD, WLS) thể hiện độ
nhạy cao hơn rõ rệt so với các mô hình có chính quy hóa (Ridge, Lasso, Elastic
Net).  Các mô hình chính quy hóa cho thấy khả năng chống chịu tốt hơn trước
nhiễu do dữ liệu thiếu, thể hiện qua các giá trị độ nhạy thấp hơn đáng kể trên
mọi chiến lược imputation.

Đối với các mô hình phi tuyến (Polynomial, RBF, Fourier), độ nhạy thường cao
hơn so với các mô hình tuyến tính, đặc biệt khi tỉ lệ dữ liệu thiếu tăng. Điều
này cho thấy các mô hình này phụ thuộc mạnh vào cấu trúc dữ liệu ban đầu, do đó
dễ bị ảnh hưởng khi dữ liệu bị thiếu và được bổ sung lại.

Tóm lại, có thể xếp hạng hiệu quả của các phương pháp bổ sung dữ liệu theo thứ
tự giảm dần: kNN, Mean, Median. Ngoài ra, khi tỉ lệ dữ liệu thiếu tăng, việc
lựa chọn mô hình (đặc biệt là các mô hình có chính quy hóa) đóng vai trò quan
trọng không kém so với việc lựa chọn phương pháp imputation.

\subsubsection{Phân lớp}

\tablename~\ref{tab:classification-sens-imputer} trình bày độ nhạy của các mô
hình phân lớp trước các tỉ lệ dữ liệu thiếu khác nhau khi áp dụng các chiến
lược bổ sung dữ liệu (mean, median).

\import{\tablespath}{classification-sens-imputer}

Trước hết, có thể nhận thấy một xu hướng nhất quán rằng khi tỉ lệ dữ liệu thiếu
tăng từ 0.1 lên 0.3, độ nhạy của tất cả các mô hình đều tăng đáng kể trên mọi
chỉ số (Accuracy, Precision macro, Recall macro và F1 macro). Điều này cho thấy
việc bổ sung dữ liệu không thể bù đắp hoàn toàn thông tin bị mất, đặc biệt khi
mức thiếu dữ liệu lớn.

So sánh giữa hai phương pháp, mean imputation thường cho kết quả tốt hơn hoặc
tương đương so với median trên hầu hết các mô hình và chỉ số. Cụ thể, các giá
trị độ nhạy khi sử dụng mean nhìn chung thấp hơn hoặc chênh lệch không đáng kể
so với median,đặc biệt ở các mô hình Softmax và LDA.

Ngược lại, median imputation có xu hướng kém ổn định hơn, thể hiện rõ ở các mô
hình nhạy như QDA. Ví dụ, tại mức thiếu 0.2 và 0.3, độ nhạy của QDA với median
tăng đột biến so với mean trên tất cả các chỉ số, cho thấy phương pháp này dễ
làm sai lệch phân phối dữ liệu trong bối cảnh hiện tại.

Xét theo từng mô hình, QDA là mô hình nhạy nhất với dữ liệu thiếu, với độ nhạy
tăng rất mạnh khi missing rate tăng, đặc biệt khi sử dụng median. Các mô hình
Softmax và LDA cho thấy độ ổn định tốt hơn, trong khi Weighted Softmax có độ
nhạy cao hơn Softmax chuẩn nhưng vẫn thấp hơn đáng kể so với QDA.

Tổng hợp lại, theo thứ tự hiệu quả giảm dần: mean, median. Ngoài ra, do hạn chế
về chi phí tính toán trên tập dữ liệu lớn, phương pháp kNN imputation không
được xem xét trong thí nghiệm này.
